% !TeX root = ../main.tex
%%%%%%%%%%%%%%%%%%%%%%%%%%%%%%%%%%%%%%%%%%%%%%%%%%%%%%%%%%%%%%%%%%%%%%
%  run mode: pdfLaTeX（兼容 XeLaTeX）
%  !Mode: "Tex:Utf-8"
%  英文摘要：定义 \abstractEn 与 \keywordsEn，由 main.tex 排版输出
%%%%%%%%%%%%%%%%%%%%%%%%%%%%%%%%%%%%%%%%%%%%%%%%%%%%%%%%%%%%%%%%%%%%%%

\abstractEn{

  With the deep integration of the Internet of Things (IoT) and
  the sixth-generation (6G) mobile communication, a huge number of
  sensing nodes continuously generate status-update data in
  real-time monitoring scenarios. Under the constraints of limited
  spectrum and energy, how to efficiently collect and schedule these
  multi-source data while keeping the system stable and reducing the
  transmission delay and staleness has become a key problem. The Age
  of Information (AoI), as an emerging metric that quantifies the
  freshness of the information held at the receiver, has attracted
  extensive attention from both academia and industry.

  This thesis studies the joint data collection and transmission
  scheduling problem in multi-source status-update scenarios for 6G
  IoT. The main contributions are threefold. First, for the network
  structure where multi-source sensors are aggregated to the base
  station through queues, a status-update and queuing-transmission
  system model is established, and the problem is formulated as a
  discrete-time optimization problem that minimizes the long-term
  average weighted AoI under the queue-stability constraint. Second,
  by combining deep learning with queuing theory, a Stacked
  Autoencoder (SAE) is exploited to compress high-dimensional state
  data in an unsupervised manner, so as to reduce the per-transmission
  load without significant information loss. Third, an online
  scheduling policy is designed based on Lyapunov drift minimization,
  which decides the optimal source to transmit in each slot according
  to the current queue state and AoI, achieving asymptotically optimal
  and stable scheduling without relying on future channel and arrival
  statistics.

  Simulation results show that, compared with classical baselines
  such as round-robin, maximum-AoI-first and maximum-queue-first
  policies, the proposed method significantly reduces the long-term
  average weighted AoI while guaranteeing queue stability, and the
  combination of the SAE compression further reduces the transmission
  bit overhead and improves the overall timeliness. This work provides
  a theoretical basis and engineering reference for large-scale data
  collection and information-timeliness guarantee in 6G IoT.

}
%这里是英文关键词
\keywordsEn{Age of Information; Stacked Autoencoder; Lyapunov Optimization;
Multi-source Scheduling; Internet of Things; 6G}
